362
6 Integration
\medskip
Thus we have shown that if the function $f$ is continuous at a point $x \in [a, b]$, then
for displacements $h$ from $x$ such that $x + h \in [a, b]$ the following equality holds:
\[
F(x + h) – F(x) = f (x)h + \alpha(h)h,
\tag{6.42}
\]
where $\alpha(h) \to 0$ as $h \to 0$.
\medskip
But this means that the function $F(x)$ is differentiable on $[a, b]$ at the point $x \in
[a, b]$ and that $F'(x) = f(x)$.
\medskip
A very important immediate corollary of Lemma 1 is the following.
\medskip
\noindent \textbf{Theorem 1} Every continuous function $f : [a, b] \to \mathbb{R}$ on the closed interval $[a, b]$
has a primitive, and every primitive of $f$ on $[a, b]$ has the form
\[
F(x) = \int_a^x f(t)\,dt + c,
\tag{6.43}
\]
where $c$ is a constant.
\medskip
Proof We have the implication $(f \in C[a, b]) \Rightarrow (f \in \mathcal{R}[a, b])$, so that by Lemma 1
the function (6.40) is a primitive for $f$ on $[a, b]$. But two primitives $F(x)$ and $\tilde{F}(x)$
of the same function on a closed interval can differ on that interval only by a con-
stant; hence $F(x) = \tilde{F}(x) + c$.
\medskip
For later applications it is convenient to broaden the concept of primitive slightly
and adopt the following definition.
\medskip
\noindent \textbf{Definition 1} A continuous function $x \to F(x)$ on an interval of the real line is
called a primitive (or generalized primitive) of the function $x \to f(x)$ defined on
the same interval if the relation $F'(x) = f(x)$ holds at all points of the interval, with
only a finite number of exceptions.
\medskip
Taking this definition into account, we can assert that the following theorem
holds.
\medskip
\noindent \textbf{Theorem 1'} A function $f : [a, b] \to \mathbb{R}$ that is defined and bounded on a closed
interval $[a, b]$ and has only a finite number of points of discontinuity has a (gen-
eralized) primitive on that interval, and any primitive of $f$ on $[a,b]$ has the form
(6.43).
\medskip
Proof Since $f$ has only a finite set of points of discontinuity, $f \in \mathcal{R}[a, b]$, and by
Lemma 1 the function (6.40) is a generalized primitive for $f$ on $[a, b]$. Here we have
taken into account, as already pointed out, the fact that by (6.41) the function (6.40)
is continuous on $[a, b]$. If $F(x)$ is another primitive of $f$ on $[a, b]$, then $F(x) - \tilde{F}(x)$
is a continuous function and constant on each of the finite number of intervals into
which the discontinuities of $f$ divide the closed interval $[a, b]$. But it then follows
from the continuity of $F(x) - \tilde{F}(x)$ on all of $[a, b]$ that $F(x) - \tilde{F}(x) = \text{const}$ on
$[a, b]$.